% Figure for Electromagnetism §B9.1 (Example: energy flow into a current-carrying resistor, after Griffiths Ex. 8.1; longitudinal section; S = E x B/mu0 points radially inward with |S| proportional to 1/s outside the wire (E uniform along the wire, B = mu0 I/(2 pi s)), B out of the page above the axis and into it below for I to the right; so the arrow lengths are drawn proportional to 1/s and the power through every coaxial cylinder is the same, VI).
\begin{tikzpicture}[>={Stealth[length=2.0mm]}, font=\small]
\fill[black!12] (0,-0.350) rectangle (6.000,0.350);
\draw[black!55, thin] (0,-0.350) rectangle (6.000,0.350);
\draw[->, thick, black!70] (-0.900,0) -- (-0.100,0) node[midway, above] {$I$};
\draw[->, thick, black!70] (6.100,0) -- (6.900,0);
\draw[->, thick, figB] (1.250,0) -- (1.950,0);
\draw[->, thick, figB] (2.650,0) -- (3.350,0);
\draw[->, thick, figB] (4.050,0) -- (4.750,0);
\node[figB, anchor=south] at (3.0,0.02) {$\vec E$};
\draw[->, thick, figG] (0.900,0.960) -- (0.900,0.540);
\draw[->, thick, figG] (2.300,0.960) -- (2.300,0.540);
\draw[->, thick, figG] (3.700,0.960) -- (3.700,0.540);
\draw[->, thick, figG] (5.100,0.960) -- (5.100,0.540);
\draw[->, thick, figG] (0.900,1.287) -- (0.900,1.013);
\draw[->, thick, figG] (2.300,1.287) -- (2.300,1.013);
\draw[->, thick, figG] (3.700,1.287) -- (3.700,1.013);
\draw[->, thick, figG] (5.100,1.287) -- (5.100,1.013);
\draw[->, thick, figG] (0.900,1.840) -- (0.900,1.660);
\draw[->, thick, figG] (2.300,1.840) -- (2.300,1.660);
\draw[->, thick, figG] (3.700,1.840) -- (3.700,1.660);
\draw[->, thick, figG] (5.100,1.840) -- (5.100,1.660);
\draw[figR, thick] (1.600,0.950) circle (0.09); \fill[figR] (1.600,0.950) circle (0.03);
\draw[figR, thick] (3.000,0.950) circle (0.09); \fill[figR] (3.000,0.950) circle (0.03);
\draw[figR, thick] (4.400,0.950) circle (0.09); \fill[figR] (4.400,0.950) circle (0.03);
\draw[->, thick, figG] (0.900,-0.960) -- (0.900,-0.540);
\draw[->, thick, figG] (2.300,-0.960) -- (2.300,-0.540);
\draw[->, thick, figG] (3.700,-0.960) -- (3.700,-0.540);
\draw[->, thick, figG] (5.100,-0.960) -- (5.100,-0.540);
\draw[->, thick, figG] (0.900,-1.287) -- (0.900,-1.013);
\draw[->, thick, figG] (2.300,-1.287) -- (2.300,-1.013);
\draw[->, thick, figG] (3.700,-1.287) -- (3.700,-1.013);
\draw[->, thick, figG] (5.100,-1.287) -- (5.100,-1.013);
\draw[->, thick, figG] (0.900,-1.840) -- (0.900,-1.660);
\draw[->, thick, figG] (2.300,-1.840) -- (2.300,-1.660);
\draw[->, thick, figG] (3.700,-1.840) -- (3.700,-1.660);
\draw[->, thick, figG] (5.100,-1.840) -- (5.100,-1.660);
\draw[figR, thick] (1.600,-0.950) circle (0.09) (1.536,-1.014) -- (1.664,-0.886) (1.536,-0.886) -- (1.664,-1.014);
\draw[figR, thick] (3.000,-0.950) circle (0.09) (2.936,-1.014) -- (3.064,-0.886) (2.936,-0.886) -- (3.064,-1.014);
\draw[figR, thick] (4.400,-0.950) circle (0.09) (4.336,-1.014) -- (4.464,-0.886) (4.336,-0.886) -- (4.464,-1.014);
\node[figR, anchor=west] at (4.550,0.95) {$\vec B$};
\node[figG, anchor=west] at (5.300,1.75) {$\vec S$};
\draw[|<->|, thin, black!55] (0.300,0) -- node[right, font=\scriptsize] {$a$} (0.300,0.350);
\draw[|<->|, thin, black!55] (0,-2.350) -- node[below, font=\scriptsize] {$L$} (6.000,-2.350);
\node[font=\scriptsize, black!70, align=left, anchor=west] at (6.350,1.200) {$|\vec S| = \dfrac{VI}{2\pi sL}$\\[2pt] power through any\\ coaxial cylinder: $VI$};
\end{tikzpicture}
